% Part: second-order-logic
% Chapter: syntax-and-semantics
% Section: semantic-notions

\documentclass[../../../include/open-logic-section]{subfiles}

\begin{document}

\olfileid{sol}{syn}{sem}
\olsection{అర్థపర భావనలు}

\begin{explain}
\emph{చెల్లుబాటుతనం}, \emph{అనుగమనం}, \emph{సంతృప్తిపరచదగినతనం}
అనే కేంద్ర తార్కిక భావనలను ద్వితీయ-స్థాయి తర్కానికి మొదటిస్థాయి
తర్కంలోలాగే నిర్వచిస్తాం; అయితే ఆధారమైన సంతృప్తి సంబంధం ఇప్పుడు
ద్వితీయ-స్థాయి \tetoken{సూత్రాలకు}{!!{formula}s} సంబంధించినది. వాస్తవానికి, ద్వితీయ-స్థాయి
\tetoken{వాక్యం}{!!{sentence}} అంటే విధేయ, ప్రమేయ చరాలతో సహా అన్ని చరాలూ బద్ధమైన
\tetoken{ఒక సూత్రమే}{!!a{formula}}.
\end{explain}

\begin{defn}[చెల్లుబాటుతనం]
ఒక వాక్యం~$!A$ \emph{చెల్లుబాటు అవుతుంది}, $\Entails !A$, అప్పుడే,
అప్పుడు మాత్రమే ప్రతి \tetoken{నిర్మాణం}{!!{structure}}~$\Struct M$కు $\Sat{M}{!A}$.
\end{defn}

\begin{defn}[అర్థపర అనుగమనం]
వాక్యాల సమితి~$\Gamma$ నుంచి వాక్యం~$!A$ \emph{అనుగమిస్తుంది},
$\Gamma \Entails !A$, అప్పుడే, అప్పుడు మాత్రమే $\Sat{M}{\Gamma}$
అయ్యే ప్రతి \tetoken{నిర్మాణం}{!!{structure}}~$\Struct M$కు $\Sat{M}{!A}$.
\end{defn}

\begin{defn}[సంతృప్తిపరచదగినతనం]
ఏదో ఒక \tetoken{నిర్మాణం}{!!{structure}}~$\Struct M$కు $\Sat{M}{\Gamma}$ అయితే,
వాక్యాల సమితి~$\Gamma$ \emph{సంతృప్తిపరచదగినది}. $\Gamma$
సంతృప్తిపరచదగినది కాకపోతే దాన్ని \emph{సంతృప్తిపరచదగనిది} అంటాం.
\end{defn}

\end{document}
